\documentclass[11pt,a4paper]{article}
\usepackage[margin=2.5cm]{geometry}
\usepackage{amsmath,amssymb}
\usepackage{booktabs}
\usepackage{parskip}
\usepackage{hyperref}
\usepackage{xcolor}
\usepackage{graphicx}

\hypersetup{colorlinks=true,linkcolor=blue,citecolor=blue,urlcolor=blue}

\newcommand{\DS}{\Delta\Sigma}
\newcommand{\gtrue}{\gamma_{\rm true}}
\newcommand{\ghat}{\hat\gamma}
\newcommand{\ceff}{c_{\rm eff}}
\newcommand{\dm}{\Delta m}
\newcommand{\Msun}{\,M_{\odot}}
\newcommand{\Mpc}{\,\mathrm{Mpc}}

\title{\textbf{201-point Gamma Scan and the Triaxial-vs-Spherical
$\DS$ Mismatch}\\[0.4em]
\large Destacking project --- combined diagnostic report}
\author{destacking project}
\date{2026 May 27}

\begin{document}
\maketitle

\begin{abstract}
We present two complementary results.  (1) A high-resolution gamma scan with
201 $\gtrue$ values from $-1$ to $+1$, $N_{\rm clusters}=5000$ per simulation,
single full-sample sampler fit per $\gtrue$ (no bootstrap), with the
single-posterior 16/84 half-width used as the error.  The weighted linear
fit recovers $\ghat = (0.809 \pm 0.069)\,\gtrue + (0.111 \pm 0.057)$,
i.e.\ a $\sim 19\,\%$ slope compression with a small additive offset.
The bias is largest in the unphysical $|\gtrue| > 0.5$ regime where the
implied per-cluster $c$ leaves the literature range $c\in[1,10]$.  (2) A
controlled spherical-vs-triaxial NFW comparison at zero shape noise shows
that the orientation-averaged triaxial signal is a $\sim 20\,\%$
\emph{uniform} downward shift relative to the spherical reference at the
same $(M,c)$ --- flat in both radius and concentration over more than a
decade of either.  A spherical NFW fit to this stacked triaxial signal at
fixed $M$ recovers $\ceff$ within a few per cent of the input $c$ (with
large statistical uncertainty), because a uniform multiplicative amplitude
offset cannot be absorbed by a single concentration parameter alone.  The
two results join up cleanly: the triaxial signal looks like a constant
\emph{amplitude} shift on the spherical model, which in the destacking
inference is forced onto $\gamma$ because the natural sink (mass via
$\log M_{\rm norm}$) is tightly priored.
\end{abstract}

\tableofcontents
\newpage

% ============================================================
\section{201-point gamma scan}
% ============================================================

\subsection{Setup}

We ran \texttt{scripts/gamma\_scan.py} with
\texttt{BOOTSTRAP=0 N\_CLUSTERS=5000 N\_GAMMA=201}.  This produces a 201-point
sweep of $\gtrue$ in $[-1,+1]$ at step $0.01$.  With bootstrap off, each
$\gtrue$ is a single sampler run on the full 5000-cluster sample; the error
bar per point is the posterior $(\gamma_{84}-\gamma_{16})/2$ from that
single fit (the previous-default $\sigma_{\rm boot} = \mathrm{std}_b(\ghat^{(b)})$
is undefined for $N_b=1$).  All other configuration follows
\texttt{gamma\_scan.py} defaults:
\texttt{R\_MIN\_FIT=0.5\,Mpc}, \texttt{F\_CEN=0.64},
\texttt{SIGMA\_OFF\_MPC=0.32/h}, tight priors
$c_0\!\sim\! N(4.57,0.05)$, $\alpha\!\sim\! N(-0.1,0.05)$,
$\beta\!\sim\! N(-0.4,0.2)$, $\gamma\!\sim\! N(0,3.5)$,
$\log M_{\rm norm}\!\sim\! N(0,0.1)$.

The run completed in $9.18\,\mathrm{h}$ wall time
(reference: \texttt{gamma\_scan\_outputs/run\_2026-05-26T18-27-50/}).

\subsection{Transfer function and linear fit}

Figure~\ref{fig:scan} shows the recovered transfer function $\ghat(\gtrue)$
and the residual.  Three overlays are added on top of the per-point data:

\begin{itemize}
  \item a weighted linear fit  $\ghat = a\,\gtrue + b$  with
        $\pm 1\sigma$ band propagated from the full covariance;
  \item a running median of $\ghat$ in non-overlapping bins of 10
        consecutive points, with the 68\,\% per-bin band;
  \item the standard error bars (single-posterior 16/84 half-width)
        on each $\gtrue$ point.
\end{itemize}

\begin{figure}[htbp]
  \centering
  \includegraphics[width=0.9\linewidth]{../gamma_scan_outputs/run_2026-05-26T18-27-50/gamma_scan_results.png}
  \caption{201-point gamma scan, single full-sample fit per $\gtrue$.  Top:
  $\ghat$ vs $\gtrue$ with the linear fit (crimson) and binned median (purple)
  overlaid.  Bottom: residuals with the same overlays in residual space.
  The linear fit gives $\ghat=(0.809\pm 0.069)\gtrue+(0.111\pm 0.057)$;
  the slope deviation from~1 is $-0.191\pm 0.069$.}
  \label{fig:scan}
\end{figure}

The numerical result is summarised in Table~\ref{tab:scan-summary}.

\begin{table}[htbp]
  \centering
  \begin{tabular}{lr}
    \toprule
    quantity & value\\
    \midrule
    linear-fit slope $a$ & $0.809 \pm 0.069$\\
    slope deviation from 1   & $-0.191 \pm 0.069$\\
    linear-fit intercept $b$ & $+0.111 \pm 0.057$\\
    mean residual            & $+0.111$\\
    RMS bias                 & $0.421$\\
    68\,\% CL on residual    & $[-0.229,\,+0.407]$\\
    mean per-point $\sigma$  & $0.141$\\
    \midrule
    \multicolumn{2}{l}{\emph{physical-regime subset: $\gtrue\in[0,0.5]$,
    $N=51$ points}}\\
    mean residual            & $+0.079$\\
    RMS bias                 & $0.383$\\
    68\,\% CL on residual    & $[-0.360,\,+0.407]$\\
    mean per-point $\sigma$  & $0.160$\\
    \bottomrule
  \end{tabular}
  \caption{201-point gamma scan summary statistics, full range and the
  physical subset $\gtrue\in[0,0.5]$.}
  \label{tab:scan-summary}
\end{table}

The fit slope of $0.81$ means the estimator is compressed toward zero by
$\sim 19\,\%$ per unit $\gtrue$ across the full range.  The intercept
$+0.111$ is a small additive offset, statistically significant at
$\sim 2\sigma$.  The same plateau-like structure we diagnosed earlier
(\texttt{2026-05-26\_gamma\_plateau\_diagnosis.pdf}) is visible: $\ghat$
saturates near $\pm 0.5$ for $|\gtrue| \gtrsim 0.6$, with a smoother
transition through the middle.

\subsection{Physical-regime caveat}

Per earlier discussion, only $\gtrue \ge 0$ is physically motivated (fossil
clusters are observed to have $\dm \ge 2$ and to be \emph{more} concentrated,
favouring $\gamma>0$).  In addition, the literature does not observe NFW
concentrations outside $c\in[1,10]$ in the cluster mass range.  With
$c_{\rm DM}\approx 4.5$ at the pivot and $\dm$ spanning $[0.01,2]$\,mag,
the implied $c_{\rm truth} = c_{\rm DM}\,(\dm/0.1)^{\gtrue}$ stays
inside $[1,10]$ for the bulk of clusters only when $\gtrue\lesssim 0.5$.
Above that, an increasing fraction of clusters has unphysical $c$, and the
bias inflates accordingly (see the bottom panel of Fig.~\ref{fig:scan}
beyond $\gtrue\!\sim\!0.5$).  The physical-regime subset in
Table~\ref{tab:scan-summary} therefore gives a fairer picture: mean residual
$+0.079$ over $\gtrue\in[0,0.5]$, with a 68\,\% CL of $[-0.36,+0.41]$
dominated by the few points at $\gtrue\to 0.5$ where the $c$-distribution
already pushes against the upper edge.

% ============================================================
\section{Spherical vs.\ triaxial NFW $\DS$ comparison}
% ============================================================

\subsection{Setup}

We isolate the triaxial-vs-spherical projection effect with a controlled
zero-shape-noise comparison
(\texttt{scripts/sphere\_vs\_triaxial.py},
output in \texttt{sphere\_vs\_triaxial\_outputs/}).  For each of $N=100$
triaxial shape draws at fixed $(z=0.4, M=10^{14}\Msun, c=5)$ we generate
one zero-shape-noise background-galaxy field and shear it twice: once with
\texttt{apply\_NFW\_shear\_region} (spherical NFW), once with
\texttt{apply\_triaxial\_NFW\_shear\_region} using the sampled shape.  The
two sheared catalogs are passed through the same
\texttt{metacal\_cluster\_lensing}\,$\to$\,\texttt{signal} pipeline used by
the gamma scan, yielding two $\DS$ profiles per cluster.  With shape noise
$=0$ and intrinsic $e_1=e_2=0$, the only signal in the recovered $\DS$ is
the lensing-induced one, so the difference between sphere and triaxial is
purely geometric projection.

The triaxial shapes are drawn from the same prior used in
\texttt{gamma\_scan.py}: $b/a\sim U(0.5,1)$, $c/a\sim U(0.3,b/a)$, isotropic
LOS.  This is a prolate-dominated population (median triaxiality
$T=0.643$, $47\,\%$ prolate, $31\,\%$ generally triaxial, $21\,\%$ oblate;
median projected $q_{\min}/q_{\max}\approx 0.68$).

\subsection{Mean profiles and the stacked fit}

Figure~\ref{fig:profiles} overlays the per-cluster triaxial profiles, the
spherical reference (population mean, identical for all clusters), the
stacked triaxial mean over orientations, and a spherical NFW model fit to
the stack with $M$ fixed and $c$ free.

\begin{figure}[htbp]
  \centering
  \includegraphics[width=0.78\linewidth]{../sphere_vs_triaxial_outputs/profiles.png}
  \caption{$\DS$ profiles at $z=0.4$, $M=10^{14}\Msun$, $c=5$, zero shape
  noise, $N=100$ shape draws.  Black: spherical NFW input (also the stacked
  spherical signal, since all clusters are identical).  Steelblue: stacked
  triaxial NFW (population mean over orientations) --- this is what a
  traditional stacking analysis would observe.  Red dashed: best-fit
  spherical NFW with $M$ fixed and $c$ free; $\ceff = 5.18 \pm 1.97$ from
  fitting only $R\ge 0.5\Mpc$, the radial range the inference pipeline uses.}
  \label{fig:profiles}
\end{figure}

Figure~\ref{fig:frac} expresses the same information as the per-cluster
fractional residual $(\DS_{\rm tri}-\DS_{\rm sph})/\DS_{\rm sph}$.  The
stacked deficit (green diamonds) is flat at $-20\,\%$ across all radial
bins.  The best-fit spherical NFW (red dotted) is within $\pm 4\,\%$ of
the input $c=5$ reference everywhere --- the fit barely moves because
varying $c$ at fixed $M$ \emph{redistributes} mass radially but cannot
produce a uniform amplitude shift.

\begin{figure}[htbp]
  \centering
  \includegraphics[width=0.78\linewidth]{../sphere_vs_triaxial_outputs/fractional_residual.png}
  \caption{Fractional residual relative to spherical NFW at $c=5$.  Shaded:
  68\,\% per-cluster band (mostly orientation-driven).  Blue squares:
  median per-cluster residual.  Green diamonds: stacked-triaxial residual,
  flat at $\sim -20\,\%$.  Red dotted: best-fit spherical NFW
  ($\ceff=5.18$) deviation from the $c=5$ reference --- only a few per cent
  because $c$ cannot fit a flat amplitude shift.}
  \label{fig:frac}
\end{figure}

\subsection{Orientation and triaxiality dependence}

Figure~\ref{fig:orient} splits the per-cluster fractional deficit by LOS
orientation.  Three roughly equal bins
($\cos\theta_{\rm LOS}<0.33$, $0.33\le\cos\theta_{\rm LOS}<0.66$,
$\cos\theta_{\rm LOS}\ge 0.66$) recover three flat-in-$R$ deficit lines at
$\sim -35\,\%$ (edge-on), $\sim -15\,\%$ (intermediate),
$\sim -11\,\%$ (face-on).  The factor of $\sim 3$ in deficit across
orientation bins is the largest single source of inter-cluster variance.

Figure~\ref{fig:tri} splits the same data by intrinsic triaxiality
$T$ (oblate, generally triaxial, prolate).  The three lines lie within
$\pm 2\,\%$ of one another --- intrinsic triaxiality is essentially absorbed
into orientation after isotropic LOS averaging.

\begin{figure}[htbp]
  \centering
  \begin{minipage}{0.49\linewidth}
    \includegraphics[width=\linewidth]{../sphere_vs_triaxial_outputs/by_orientation.png}
    \caption{Deficit split by LOS orientation.}
    \label{fig:orient}
  \end{minipage}\hfill
  \begin{minipage}{0.49\linewidth}
    \includegraphics[width=\linewidth]{../sphere_vs_triaxial_outputs/by_triaxiality.png}
    \caption{Deficit split by intrinsic triaxiality $T$.}
    \label{fig:tri}
  \end{minipage}
\end{figure}

\subsection{$c$-dependence: the deficit is essentially flat in $c$}

We repeated the same comparison at $N=50$ shape draws per concentration,
sweeping $c\in[0.5,30]$ in 13 log-spaced points
(\texttt{scripts/deficit\_vs\_concentration.py}).  Figure~\ref{fig:deficitc}
(top) shows the mean and median per-cluster fractional deficit averaged
over $R\ge 0.5\Mpc$ as a function of $c$.  Across more than a $60\times$
range in $c$, the mean deficit varies by less than $1$ percentage point
($-20.6\,\%$ at $c=0.5$ to $-21.4\,\%$ at $c=5$).  The orientation splits
(edge-on, intermediate, face-on) are also flat in $c$.

\begin{figure}[htbp]
  \centering
  \includegraphics[width=0.78\linewidth]{../sphere_vs_triaxial_outputs/deficit_vs_concentration.png}
  \caption{Triaxial-vs-spherical deficit vs.\ concentration (top) and
  stacking-inferred $\ceff$ vs.\ input $c$ (bottom).  The deficit is
  essentially constant in $c$.  The stacking fit recovers $\ceff$ within
  $\sim 5$--$10\,\%$ of the input $c$ across the physical range $c\in[1,10]$
  (green band), with large statistical uncertainty.}
  \label{fig:deficitc}
\end{figure}

The bottom panel shows the effective concentration $\ceff$ obtained by
fitting a spherical NFW (with $M$ fixed) to the stacked triaxial signal at
each input $c$.  $\ceff$ tracks $c_{\rm input}$ with a $+5$ to $+10\,\%$
positive bias inside the physical range $[1,10]$, growing to
$+21\,\%$ at $c_{\rm input}=30$ where the profile is so peaked that the
$R\ge 0.5\Mpc$ window barely catches the relevant features.  The error
bars are large ($\sim 30$--$100\,\%$) because a single $c$ parameter cannot
reproduce a uniform amplitude shift: the fit is essentially degenerate, and
the bias would be much better absorbed by adjusting $M$.

Table~\ref{tab:ceff} gives the numbers.

\begin{table}[htbp]
  \centering
  \begin{tabular}{rrrr}
    \toprule
    $c_{\rm input}$ & mean deficit [\%] & $\ceff$ & frac.\ bias\\
    \midrule
    $0.50$  & $-20.6$ & $0.54\pm 0.16$   & $+7.5\,\%$\\
    $0.99$  & $-21.0$ & $1.05\pm 0.29$   & $+6.3\,\%$\\
    $1.96$  & $-21.3$ & $2.06\pm 0.62$   & $+5.3\,\%$\\
    $3.87$  & $-21.4$ & $4.08\pm 1.52$   & $+5.3\,\%$\\
    $5.45$  & $-21.4$ & $5.77\pm 2.49$   & $+5.9\,\%$\\
    $10.78$ & $-21.4$ & $11.72\pm 7.07$  & $+8.7\,\%$\\
    $21.33$ & $-21.3$ & $24.62\pm 21.5$  & $+15.5\,\%$\\
    $30.00$ & $-21.2$ & $36.34\pm 38.5$  & $+21.1\,\%$\\
    \bottomrule
  \end{tabular}
  \caption{Stacking-fit $\ceff$ vs.\ input $c$.  Bias stays at
  $5$--$10\,\%$ across the physical $c\in[1,10]$ range, with very large
  statistical uncertainty.}
  \label{tab:ceff}
\end{table}

% ============================================================
\section{Joining the two: how triaxiality drives the gamma scan bias}
% ============================================================

The triaxial result is a near-uniform $\sim -20\,\%$ amplitude offset
on the spherical NFW prediction.  Three facts decide how that offset
propagates to the inference:

\begin{enumerate}
\item \textbf{The offset is flat in $R$.}  No reshuffling of the radial
      profile from changing concentration can undo it.  This is why the
      best-fit-$c$ stacking analysis shows almost no $c$ bias even though
      the data is $20\,\%$ low.

\item \textbf{The offset is flat in $c$.}  No reshuffling of the cluster
      $c$-distribution (i.e.\ no choice of $\gamma$ that changes the
      $c$-distribution shape via the $\dm$ dependence) can undo it either.
      A naive multiplicative correction would soak up the offset, but
      \texttt{test\_amplitude\_boost.py} confirmed that this does not move
      the gamma-scan peak --- the OLS-truth coefficients in the inference
      already absorb the population-mean amplitude implicitly via their
      $c$ values, so $A_{\rm opt}\approx 1$.

\item \textbf{The natural sink is mass.}  A uniform amplitude offset on
      $\DS \propto M$ is naturally absorbed by lowering $M$.  But the
      inference puts $\log M_{\rm norm}\sim N(0,0.1)$, which limits the
      mass adjustment to $\pm 10\,\%$ in $\log_{10}M_{\rm halo}$
      ($\pm 26\,\%$ in linear $M$).  $\gamma$ has prior width $3.5$ and
      ample room to absorb extra variance.
\end{enumerate}

Putting these together: the triaxial-induced $\sim -20\,\%$ amplitude
offset is partially absorbed by a few per cent $\log M_{\rm norm}$ shift,
and the remainder is forced onto $\gamma$ via the per-cluster scatter
($-3\,\%$ to $-35\,\%$ depending on LOS orientation).  Because the
orientation distribution is independent of $\dm$ in the simulation, the
per-cluster residuals don't correlate with $\dm$ in a way $\gamma$ can
neatly absorb; the result is a compressed transfer function with slope
$\sim 0.8$ and a small positive offset, exactly what
Fig.~\ref{fig:scan} measures.

The $\sim 19\,\%$ slope compression of the 201-point gamma scan is therefore
the same physical effect as the $\sim 20\,\%$ flat amplitude deficit measured
in the sphere-vs-triaxial test --- just funnelled through the tight
inference priors.

% ============================================================
\section{Practical implications}
% ============================================================

\begin{itemize}
\item In the physical regime $\gtrue\in[0,0.5]$ where the simulated cluster
      sample is realistic, the inference has a residual bias of
      $\mathrm{mean}=+0.08$ with 68\,\% CL $[-0.36,+0.41]$ over 51 points.
      This is comparable to the per-point sampler $\sigma$ ($\sim 0.16$),
      so individual $\ghat$ values cannot be trusted at the
      $\sigma_\gamma \lesssim 0.1$ level the project ultimately needs without
      a forward-model correction.

\item The bias is set by the spherical-vs-triaxial mismatch, which is a
      geometric projection effect tied to the cluster-sample shape
      distribution.  Real survey data carry no observable for per-cluster
      LOS orientation, so this bias cannot be removed cluster-by-cluster
      in the inference.  Calibration via simulations (computing the
      $\ghat(\gtrue)$ transfer function and inverting it) is the only
      avenue inside the project scope.

\item For the canonical calibration value $\gtrue=0.2$, the previous
      bootstrap scan gave $\ghat\approx 0.45$ (bias $+0.25$); the
      201-point scan in this neighbourhood is consistent
      (Fig.~\ref{fig:scan}, bottom panel).  A multiplicative slope-and-offset
      correction $\gamma_{\rm corr} = (\ghat - b)/a$ with the measured
      $a,b$ would in principle remove the linear part of the bias, leaving
      a residual at the level of the per-point scatter ($\sim 0.1$--$0.2$).
\end{itemize}

% ============================================================
\section*{Files}
% ============================================================

\begin{itemize}
\item \emph{Data:}
  \texttt{gamma\_scan\_outputs/run\_2026-05-26T18-27-50/} --- 201-point scan.\\
  \texttt{sphere\_vs\_triaxial\_outputs/} --- controlled comparison data.
\item \emph{Plots referenced in this report:}
  \texttt{gamma\_scan\_results.png},
  \texttt{profiles.png},
  \texttt{fractional\_residual.png},
  \texttt{by\_orientation.png},
  \texttt{by\_triaxiality.png},
  \texttt{deficit\_vs\_concentration.png}.
\item \emph{Scripts:}
  \texttt{scripts/gamma\_scan.py} (with new \texttt{BOOTSTRAP},
  \texttt{N\_GAMMA} env vars and single-fit $\sigma$ fix);
  \texttt{scripts/replot\_gamma\_scan\_results.py} (linear-fit and
  binned-median overlays);
  \texttt{scripts/sphere\_vs\_triaxial.py},
  \texttt{scripts/deficit\_vs\_concentration.py},
  \texttt{scripts/add\_stacked\_comparison.py},
  \texttt{scripts/plot\_triaxiality\_distribution.py}.
\item \emph{Related earlier report:}
  \texttt{2026-05-26\_gamma\_plateau\_diagnosis.pdf} --- the
  spherical-vs-triaxial mismatch was first diagnosed there.
\item \emph{Related earlier report:}
  \texttt{2026-05-26\_sphere\_vs\_triaxial.pdf} --- the first controlled
  sphere-vs-triaxial $\DS$ measurement at $c=5$, before the $c$-dependence
  scan and stacking-fit additions in this report.
\end{itemize}

\end{document}
